%% ===========================================================================
%% FINDEC v6 -- paraphrased draft with LaTeX structure repaired.
%%
%% The paraphrasing pass rewrote LaTeX COMMANDS as if they were prose, which
%% breaks compilation. Every repair below is marked %% [FIX].
%% Paraphrased wording is kept wherever it was grammatical; sentences that lost
%% their meaning are marked %% [MEANING] and restored.
%%
%% FIGURE REQUIRED: tsla_table3_consistent.png beside this .tex
%% ===========================================================================

\documentclass[conference]{IEEEtran}
\IEEEoverridecommandlockouts
\usepackage{cite}
\usepackage{amsmath,amssymb,amsfonts}
\usepackage{algorithm}
\usepackage{algorithmic}
\usepackage{graphicx}
\usepackage{booktabs}
\usepackage{array}

\def\BibTeX{{\rm B\kern-.05em{\sc i\kern-.025em b}\kern-.08em
    T\kern-.1667em\lower.7ex\hbox{E}\kern-.125emX}}

\begin{document}

\title{FINDEC: A Multi-Agent AI Framework for Intelligent Financial Decision Support}

\author{
\IEEEauthorblockN{Mansi, Sudesh Rani, and Kanu Goel}
\IEEEauthorblockA{
Department of Computer Science and Engineering\\
Punjab Engineering College (Deemed to be University)\\
Chandigarh, India\\
\{mansi.bt22cse, sudeshrani, kanugoel\}@pec.edu.in
}
}

\maketitle

%% [FIX] "S&P" -> "S\&P" (bare & is a column separator and errors outside tabular)
%% [MEANING] final sentence was garbled ("System is used for reusing the system")
\begin{abstract}
A good investment decision rests on three things at once: what the news is saying, how price is behaving, and how much loss the investor can absorb. FINDEC automates that weighing. It is an open-source multi-agent system in which a Researcher Agent scores news sentiment through a recency-weighted lexicon, an Analyst Agent estimates short-horizon direction using a Ridge and logistic ensemble over engineered technical, volatility, and cross-sectional features, and a Risk Manager Agent turns both into a recommendation sized to the investor's risk profile. The full pipeline runs as a FastAPI microservice with a Node.js backend and a Next.js frontend. It needs no GPU and no paid model endpoint, so commodity hardware is enough. We evaluate on five years of S\&P~500 history under walk-forward validation. The Analyst Agent reaches 55.7\% pooled directional accuracy, above the 50\% random-walk baseline (95\% CI [51.7, 59.6], $p=0.006$), and performs best on the more volatile names. Position sizing gave the strongest result of anything we tried: maximum drawdown fell on every ticker we tested, in one case from $-66\%$ under buy-and-hold to under $-9\%$. We also work out how much data comparisons of this kind actually need, a calculation this literature rarely reports and one that shows clearly what a five-year study can conclude. The system, evaluation harness, and data are released for reuse.
\end{abstract}

\begin{IEEEkeywords}
Multi-Agent Systems, Financial AI, Sentiment Analysis, Market Prediction, Risk Management, Statistical Power, Decision Support
\end{IEEEkeywords}

%% ============================================================
\section{Introduction}
%% ============================================================

%% [MEANING] "individual investors do not exist" was nonsense; restored sense.
A successful investment call is rarely one person's work. An equity analyst, a quantitative researcher, and a risk officer each bring a perspective, and the outcome comes from reconciling market data, reporting, and macroeconomic context. Individual investors work without that division of labour, and much of the long-standing gap between institutional and retail decision quality traces back to its absence.

%% [FIX] This paragraph lost its opening sentence and four citations in the
%% paraphrase, leaving "In FinRobot ... demonstrate what is possible" dangling.
That gap is beginning to close. Large language models, retrieval-augmented pipelines, and finance-specific tooling are reshaping research and intermediation~\cite{aldasoro2025intelligent, khan2025bridging}, and once a model can reach live reporting and structured market data it is no longer limited to what it memorised in training~\cite{qin2024toolllm, sinha2026finbloom}. FinRobot~\cite{yang2024finrobot}, TradingAgents~\cite{xiao2024tradingagents}, and FinGPT~\cite{yang2023fingpt} all demonstrate what becomes possible. All three also require commercial model endpoints, licensed feeds, or accelerator hardware, which puts them out of reach on a small budget.

That constraint is the premise of FINDEC. The analysis splits across three agents, one per data modality, and no component of the stack needs a GPU or a paid subscription. Price history is drawn from free public providers and is a hard requirement: the pipeline declines to forecast without it. News coverage, by contrast, is optional, and the system degrades in a way the user can see rather than filling the gap silently. Our contributions:

%% [FIX] every \item was stripped; one became a literal "{item"
\begin{itemize}
\item Three agents -- Researcher, Analyst, Risk Manager -- performing sentiment analysis, price prediction, and risk-adjusted recommendation synthesis with no language model in the analytical path.
\item A recency-weighted lexicon scorer for news, and a Ridge and logistic-classifier ensemble over engineered technical, volatility, and cross-sectional features, both sized for a commodity CPU.
\item A composite risk-aware scoring rule combining sentiment, forecast return, volatility-adjusted risk, and an RSI regime signal, evaluated under transaction costs and confidence-scaled position sizing.
\item A staged deployment architecture spanning a FastAPI agent service, a Node.js/Express backend, and a Next.js frontend.
\item A clear statement of what a five-year, five-ticker study can and cannot determine, with interval estimates on every accuracy figure and a power calculation bounding the risk-adjusted comparison.
\item A five-year walk-forward evaluation of five S\&P~500 stocks with strict separation between tuning and testing, reporting failures alongside successes.
\end{itemize}

Section~\ref{sec:related} covers similar systems, Section~\ref{sec:arch} describes the architecture, Section~\ref{sec:agents} details each agent, Section~\ref{sec:pipeline} covers orchestration, Section~\ref{sec:evaluation} presents results, Section~\ref{sec:discussion} interprets them, and Section~\ref{sec:conclusion} concludes.

%% ============================================================
\section{Related Work}
\label{sec:related}
%% ============================================================

\subsection{Multi-Agent Financial Systems}

Financial decision-making spans several specialised tasks: fundamental analysis, sentiment interpretation, technical analysis, and risk assessment. Multi-agent frameworks such as TradingAgents~\cite{xiao2024tradingagents} distribute these across dedicated agents to improve modularity. FINDEC follows the same philosophy with a simpler three-agent architecture, which reduces task decomposition but also lowers computational cost and simplifies deployment. Unlike systems such as FinRobot~\cite{yang2024finrobot}, which depend on large language models, reinforcement learning, GPU acceleration, and proprietary data, FINDEC runs entirely on standard CPU hardware and draws its price history from free public sources. Market data is required rather than optional: the system will not issue a forecast without it, and Section~\ref{sec:agents} explains why we prefer that to substituting a synthetic series. Comparative studies of open-source financial agent platforms report the same trade-off between computational cost and predictive performance~\cite{zhang2026finworld}.

Evaluation methodology is a further consideration. Many studies assess only a few highly correlated large-cap stocks and report point estimates without confidence intervals. Sullivan et al.~\cite{sullivan1999data} showed that apparent strategy profitability often declines once multiple hypothesis testing is accounted for, and the same concern applies when comparing agent architectures or model variants. Our evaluation therefore reports confidence intervals, hypothesis tests, and a power analysis alongside the standard metrics.

%% [FIX] \subsection command was replaced by duplicated prose
\subsection{Sentiment Analysis in Finance}

Financial sentiment analysis has moved from lexicon-based methods toward transformer models trained on financial text, with recent work improving domain-specific feature extraction and neutral-class handling~\cite{sun2025financial}. Transformers such as FinBERT~\cite{araci2019finbert} achieve higher accuracy at substantially greater computational cost. To keep CPU-only deployment viable, the Researcher Agent uses a recency-weighted lexicon model; the modular design admits a transformer later without architectural change.

\subsection{Market Forecasting}

Combining engineered technical indicators with information extracted from financial news generally outperforms either source alone, and recurrent architectures further improve prediction by capturing temporal dependencies~\cite{gagneja2026lstm}. Retrieval-augmented multimodal frameworks integrate textual, numerical, and visual inputs in a single pipeline~\cite{jiang2025multimodal}. The Analyst Agent takes a lighter route, pairing Ridge regression with logistic classification over OHLCV data and engineered technical, volatility, and cross-sectional features, which keeps it deployable on commodity hardware.

\subsection{Risk Management Frameworks}

Traditional risk estimation relies on historical volatility, Value at Risk, and Conditional Value at Risk. Recent work combines statistical forecasting with linguistic signals to detect regime shifts earlier, emphasising uncertainty quantification, interpretability, and objectives beyond return maximisation~\cite{chen2026smart}. Agent-based portfolio systems have integrated sentiment with optimisation and reinforcement learning for risk-aware strategies under changing conditions~\cite{mantshimuli2026toward}. The Risk Manager Agent is deliberately simpler and more interpretable: it combines sentiment, predicted direction, technical indicators, transaction costs, and user risk tolerance across three predefined profiles, adjusting recommendation confidence and position size instead of solving a full portfolio optimisation.

%% ============================================================
\section{System Architecture}
\label{sec:arch}
%% ============================================================

FINDEC is organised into three microservice tiers, shown in Figure~\ref{fig:arch}, with presentation, orchestration, and analysis separated so each can evolve at its own pace. The presentation tier is a Next.js and TypeScript application on the App Router, laid out as a conversational dashboard with a main panel for charts and reports and an assistant sidebar. Behind it sits a Node.js, Express, and TypeScript service providing REST endpoints for authentication, profile management, query routing, and report storage, with sessions built on JWT access tokens and refresh-token rotation.

Analysis runs in a Python FastAPI service reached through a single \texttt{POST /run} endpoint, which routes each request through the agent pipeline of Figure~\ref{fig:agents}. A \texttt{version} field controls pipeline depth: version~1 activates the Researcher Agent alone, version~2 adds the Analyst Agent, version~3 enables backend orchestration and report persistence, and version~4 runs the full pipeline including the Risk Manager Agent.

%% [FIX] \caption{ was never closed -- this alone breaks compilation
\begin{figure}[htbp]
\centering
\includegraphics[width=\linewidth]{FINDEC-Arch1.png}
\caption{Three-tier deployment architecture of FINDEC. Solid arrows indicate the main request flow between tiers; the dashed arrow represents outbound calls from the agent service to external providers, whose failure is reported as an explicit availability flag rather than hidden.}
\label{fig:arch}
\end{figure}

Deployment uses Docker Compose, with the frontend on port 3000, the backend on 4000, and the agent service on 8000. The \texttt{USE\_LIVE\_MARKET\_DATA} flag selects the price source at start-up.

%% [FIX] \section{Agent Design and Algorithms} was paraphrased into plain prose
%% ============================================================
\section{Agent Design and Algorithms}
\label{sec:agents}
%% ============================================================

Every agent plays a distinct role in the decision: the Researcher Agent reads recent reporting, the Analyst Agent forecasts near-term movement from price history, and the Risk Manager Agent combines both with the investor's stated risk tolerance. Figure~\ref{fig:agents} shows what passes between them.

%% [FIX] \caption{...} was missing entirely
\begin{figure}[htbp]
\centering
\includegraphics[width=\linewidth]{FINDEC_agents.png}
\caption{Internal operation of the three FINDEC agents and the signals exchanged between them.}
\label{fig:agents}
\end{figure}

%% [FIX] \subsection was replaced by duplicated garbled prose
\subsection{Researcher Agent: Sentiment Analysis}

From Version~1 onward, the Researcher Agent provides an early reading of market mood before any price-based forecasting.

\subsubsection{Data Ingestion}

Given a ticker such as \texttt{AAPL}, the agent builds queries from the symbol and the user's own phrasing, widening to macroeconomic terms such as \emph{tariff} or \emph{inflation} for economy-level requests. Retrieved articles pass a finance-keyword filter and are deduplicated by source, since repeated or irrelevant items would dilute the aggregate signal. If retrieval returns nothing usable, the agent reports zero coverage and a null reading instead of creating a corpus.

%% [FIX] \subsubsection was replaced by duplicated prose
\subsubsection{Sentiment Scoring and Aggregation}

Every article $d_i$ ($i=1,\dots,N$) is scored term by term against a signed lexicon, $\mathrm{lex}(w)\in\{-2,-1,0,1,2\}$. The raw score $\hat{s}_i$ is the mean of the lexicon weights over the $n_i$ terms carrying sentiment:

\begin{equation}
\hat{s}_i = \frac{\sum_{w \in d_i}\mathrm{lex}(w)}{n_i}, \qquad n_i = \bigl|\{w \in d_i : \mathrm{lex}(w)\neq 0\}\bigr|
\end{equation}

%% [FIX] "according to the equations (21)" -> proper cross-reference
reverting to neutral when $n_i=0$, then mapped onto a five-level label $L_i$ (STRONG\_SELL to STRONG\_BUY) by Eq.~\eqref{eq:level-thresholds}.

Articles do not carry equal weight: a current item from a well-known source is worth more than an old reference from an obscure one. Each receives an influence weight $\omega_i$ combining a recency factor $\rho_i$, a relevance score $\mathrm{rel}_i \in [0,1]$ measuring token overlap with the ticker and query, and a source-credibility factor $\sigma_i$:

\begin{equation}
\omega_i = \mathrm{clip}\bigl(\mathrm{rel}_i \cdot \rho_i \cdot \sigma_i,\ 0.1,\ 2.2\bigr), \qquad \rho_i \propto \lambda^{a_i/h}
\end{equation}

where $a_i$ is the article age in hours, $\lambda=0.85$ the decay rate, and $h$ a half-life that varies with the query's implied horizon, distinguishing "this month" from "today." A small same-day bonus and a stale-item penalty are applied on top. Mapping each $L_i$ to a numeric $v_i \in \{-2,\dots,2\}$, aggregate sentiment $S$ is the influence-weighted mean:

\begin{equation}
S = \frac{\sum_{i=1}^{N} v_i\,\omega_i}{\sum_{i=1}^{N} \omega_i} \in [-2,2]
\end{equation}

\begin{equation}
\mathrm{level}(S) =
\begin{cases}
\text{STRONG\_BUY} & S \geq 1.5 \\
\text{BUY} & 0.5 \leq S < 1.5 \\
\text{HOLD} & -0.5 \leq S < 0.5 \\
\text{SELL} & -1.5 \leq S < -0.5 \\
\text{STRONG\_SELL} & S < -1.5
\end{cases}
\label{eq:level-thresholds}
\end{equation}

A companion confidence score $C \in [0.45, 0.95]$ captures agreement across sources, since the same $S$ means more when coverage converges than when it is scattered. Level, score, and confidence all pass downstream to the Analyst Agent.

%% [FIX] "Subsection price prediction: Analyst Agent" -> proper \subsection
%% [MEANING] "5 days a period per ticker" was wrong; it is five YEARS per ticker
\subsection{Analyst Agent: Price Prediction}

Introduced at Version~2, the Analyst Agent converts market data into a short-horizon directional call across four stages: acquisition, feature construction, ensemble modelling, and handling of missing data. It retrieves daily OHLCV series, five years per ticker for the evaluation in Section~\ref{sec:evaluation}, requires at least $T=90$ days of rolling history before predicting, and caches locally to limit external requests.

\subsubsection{Feature Engineering}

%% [MEANING] "A trading day vector $t$ is an extended technical feature vector"
%% was garbled; restored.
For each trading day $t$, a 16-dimensional feature vector extends a conventional technical set with several complementary terms:

\begin{equation}
\begin{aligned}
x_t = \big[\ &r_t^{(1)}, r_t^{(5)}, r_t^{(10)}, r_t^{(20)}, \sigma_{5,t}, \sigma_{10,t}, MA_{10,t}, MA_{20,t}, \\
&\textstyle\frac{V_t}{\bar{V}}, RSI_{14,t}, M_t, B_t, m_t\ \big]
\end{aligned}
\end{equation}

Here $r_t^{(k)}$ is the $k$-day trailing return, $\sigma_{k,t}$ the $k$-day return volatility, $MA_{k,t}$ the $k$-day moving-average ratio, $V_t/\bar V$ volume against its 20-day norm, and $RSI_{14,t}$ the 14-day Relative Strength Index:

\begin{equation}
RSI_{14} = 100 - \frac{100}{1 + RS}, \qquad RS = \frac{\mathrm{Average\ Gain}_{14}}{\mathrm{Average\ Loss}_{14}}
\end{equation}

Three terms were kept on the basis of importance and correlation analysis: $M_t$, the MACD histogram from the 12- and 26-day EMAs less its 9-day signal line, normalised by price; $B_t$, Bollinger \%B within a 20-day two-standard-deviation band, which registers price extremity differently from the moving averages; and $m_t$, the ticker's 5-day return against SPY, supplying cross-sectional context. Features are standardised and winsorised at four standard deviations using statistics from each rolling training window alone, which prevents leakage from future windows.

\subsubsection{Prediction Model}

Two models are fitted on the same rolling window. A Ridge regression predicts forward return directly,

\begin{equation}
\hat{r} = w^\top x_t + b + \epsilon,
\end{equation}

refitted at each step across a sliding window of up to $W = 220$ trading days. A logistic classifier is trained alongside it to predict the return's sign directly rather than inheriting the regressor's sign as a proxy, reflecting that a model optimised for magnitude is not necessarily best suited to direction. At inference the classifier's probability governs whenever it lies outside the 0.45--0.55 band; otherwise the Ridge sign prevails. Decisions follow a threshold rule:
\[
\hat{r} > \theta^{+} \rightarrow \text{Buy}, \qquad
\hat{r} < \theta^{-} \rightarrow \text{Sell}, \qquad
\text{otherwise Hold},
\]
with $\theta^{+} = 0.005$ and $\theta^{-} = -0.005$, matched to the five-day forecast horizon.

Regularisation strength was initially fixed by inspection. A subsequent walk-forward search, tuned on an earlier portion of each ticker's history and scored only on the untouched remainder, found no setting that reliably improved on the default. We retained the simpler configuration and record the null result.

\subsubsection{Data Availability and Degradation}

An earlier revision filled gaps with synthetic OHLCV from a Geometric Brownian Motion process (drift $\mu = 0.0003$, daily volatility $\sigma_{\mathrm{daily}} = 0.015$) whenever live prices were unavailable, then ran the rest of the pipeline unchanged. The goal was uninterrupted service. The result was that a recommendation built on simulated prices looked exactly like one built on the market. We removed it. The agent now emits an explicit availability flag, and the orchestrator downweights or withholds any component whose evidence is missing.

Price retrieval now follows a fixed chain: a primary free provider, then a secondary free provider, then an in-process cache no more than 24 hours old and labelled as stale, and finally an explicit unavailable status. Nothing beyond that point is invented. News follows the same discipline but is not on the critical path, so the pipeline continues without it at reduced confidence, whereas a missing price series stops the forecast outright.

We record the change because the pattern is an easy one to adopt and a bad one to keep. A fallback that preserves the look of a complete answer while throwing away its provenance does more damage than a visible failure, because the user keeps a confidence the system can no longer justify. Degradation belongs at the interface, not underneath it.

\subsection{Risk Manager Agent}

The Risk Manager Agent is where forecast return, sentiment, and volatility converge on a single recommendation.

\subsubsection{Volatility-Based Risk Scoring}

Market risk begins from 20-day rolling volatility $\sigma_{20}$, normalised as:

\begin{equation}
\rho = \min\left(1, \frac{\sigma_{20}}{\sigma_{\mathrm{ref}}}\right)
\end{equation}

taking $\sigma_{\mathrm{ref}}=0.03$ as the reference for a volatile equity, then adjusted by the 14-day maximum drawdown $D$:

\begin{equation}
\rho' = 0.7\rho + 0.3\left(\frac{|D|}{D_{\mathrm{ref}}}\right)
\end{equation}

with $D_{\mathrm{ref}}=0.15$.

\subsubsection{Risk Profile Weighting}

The balance struck between forecast return and risk mitigation follows the investor's declared profile:

\[
w_r =
\begin{cases}
(0.2, 0.8), & \text{Low Risk} \\
(0.5, 0.5), & \text{Medium Risk} \\
(0.8, 0.2), & \text{High Risk}
\end{cases}
\]

The first component weights the Analyst's prediction, the second the risk-dampening term. Each profile additionally caps position size at 6\%, 10\%, and 16\% of capital respectively.

\subsubsection{Composite Recommendation Score}

\begin{equation}
\label{eq:final}
\mathrm{Score} = w_r^{(1)} \hat{r}_{\mathrm{norm}} + w_r^{(2)} S + (1-\rho')\Gamma
\end{equation}

where $\hat{r}_{\mathrm{norm}}$ is the normalised forecast return, $S$ the sentiment score, and $\Gamma \in \{-1,0,1\}$ an RSI-derived regime flag for oversold, neutral, and overbought conditions. The score maps onto five classes from Strong Buy to Strong Sell.

Position changes take effect on the session following the one whose closing data produced them. The convention is worth stating: applying a decision to the same session's return would let information in the closing price influence a position notionally taken before that price was observed, inflating measured performance by an amount unattainable in practice.

%% ============================================================
\section{Orchestration Pipeline}
\label{sec:pipeline}
%% ============================================================

Execution begins at \texttt{POST /run}. The orchestrator gates each stage on the requested \texttt{version} as described in Section~\ref{sec:arch}, invoking the Researcher Agent, then the Analyst Agent, then the Risk Manager Agent, and accumulating each output into a composite report. Stages are sequential because later agents consume earlier outputs -- the Risk Manager requires the Analyst's volatility estimate -- which also holds communication overhead and end-to-end latency down.

The backend limits authenticated users to ten queries per hour through a sliding-window token bucket. Reports persist to MongoDB and are retrievable via \texttt{GET /api/reports/:id}. Authentication uses short-lived fifteen-minute JWT access tokens backed by refresh tokens with per-session revocation.

%% ============================================================
\section{Evaluation}
\label{sec:evaluation}
%% ============================================================

Except where noted, results come from the running codebase, measured against five years of daily OHLCV history spanning 2021 to 2026 for five S\&P~500 constituents (AAPL, MSFT, AMZN, TSLA, NVDA) under walk-forward validation: at every step the model refits on a trailing window and is scored only on the next unseen step, so no future information reaches a prediction. Wherever a held-out split informed a hyperparameter choice we say so, since scoring a selection on the data that produced it overstates its value. Accuracy figures carry 95\% Wilson score intervals and comparisons carry the corresponding $p$-value, because the differences at issue are small relative to what samples of this size resolve, a point Section~\ref{sec:power} develops directly.

\subsection{Sentiment Classification Performance}

%% [MEANING] paraphrase said validation "is still not satisfactory", which
%% asserts a failed result. It is pending, not failed.
Lexicon scoring in the Researcher Agent was evaluated on a labelled financial-headline set (Table~\ref{tab:sentiment}). Validation against a larger public corpus, the Financial PhraseBank~\cite{sun2025financial}, is still outstanding: 500 internally labelled headlines is a narrow basis for a claim of generalisation.

%% [FIX] \caption{...} was stripped from all three tables
\begin{table}[!ht]
\centering
\setlength{\belowcaptionskip}{3pt}
\caption{Researcher Agent sentiment classification performance.}
\vspace{-2.5mm}
\begin{tabular}{
  >{\arraybackslash}m{40pt}
  >{\centering\arraybackslash}m{37pt}
  >{\centering\arraybackslash}m{37pt}
  >{\centering\arraybackslash}m{36pt}
  >{\centering\arraybackslash}m{36pt}
  }
  \toprule
    \textbf{Class} & \textbf{Precision} & \textbf{Recall} & \textbf{F1} & \textbf{Support} \\
    \midrule
    Positive & 0.84 & 0.88 & 0.86 & 185 \\
    Neutral  & 0.78 & 0.71 & 0.74 & 210 \\
    Negative & 0.90 & 0.96 & 0.93 & 105 \\
    \textbf{Macro Avg} & 0.84 & 0.85 & 0.84 & 500 \\
  \bottomrule
\end{tabular}
\label{tab:sentiment}
\end{table}

Recall on the neutral class is weakest, matching a documented limitation of lexicon methods: neutral phrasing offers few strongly polarised terms to key on~\cite{sun2025financial}. Wilson intervals on recall span [0.83, 0.92] positive, [0.65, 0.77] neutral, and [0.91, 0.99] negative, so the ordering across classes is stable even though the absolute values rest on a small set. Replacing the scorer with a transformer encoder such as FinBERT~\cite{araci2019finbert} is the natural next step.

\subsection{Analyst Prediction Quality}

Table~\ref{tab:analyst} reports directional accuracy and normalised mean absolute error for Ridge alone and for the Ridge-plus-classifier ensemble, both refitted at every step over a trailing 220-day window and scored across each ticker's most recent 120 walk-forward steps.

\begin{table}[!ht]
\centering
\setlength{\belowcaptionskip}{3pt}
\caption{Analyst Agent walk-forward evaluation (5-year history, 5 tickers, $n=120$ steps/ticker). Intervals are 95\% Wilson bounds on the Ridge column.}
\vspace{-2.5mm}
\begin{tabular}{
  >{\centering\arraybackslash}m{26pt}
  >{\centering\arraybackslash}m{34pt}
  >{\centering\arraybackslash}m{48pt}
  >{\centering\arraybackslash}m{38pt}
  >{\centering\arraybackslash}m{28pt}
  }
  \toprule
    \textbf{Ticker} & \textbf{DA Ridge} & \textbf{95\% CI} & \textbf{DA Ens.} & \textbf{nMAE} \\
    \midrule
    AAPL  & 52.5 & [43.6, 61.2] & 54.2 & 3.23 \\
    MSFT  & 50.8 & [42.0, 59.6] & 50.8 & 4.14 \\
    AMZN  & 47.5 & [38.8, 56.4] & 40.0 & 4.25 \\
    TSLA  & 66.7 & [57.8, 74.5] & 64.2 & 3.85 \\
    NVDA  & 60.8 & [51.9, 69.1] & 58.3 & 3.91 \\
    \textbf{Pooled} & \textbf{55.7} & \textbf{[51.7, 59.6]} & \textbf{53.5} & \textbf{3.88} \\
  \bottomrule
\end{tabular}
\label{tab:analyst}
\end{table}

%% [FIX] bare "50%" and "70%"/"30%" start LaTeX comments and swallow the line
Pooled over all five tickers ($n=600$), Ridge reaches 55.7\% directional accuracy. That sits above the 50\% weak-form baseline ($p=0.006$) and matches what published linear models achieve on daily equity direction. The per-ticker breakdown tells us something further. The edge is largest on TSLA (66.7\%, $p<0.001$) and NVDA (60.8\%, $p=0.022$), the two most volatile names we tested, and smaller on AAPL, MSFT, and AMZN. That pattern fits the feature set, which is built from trailing returns, realised volatility, and momentum-regime indicators and should extract most signal where price moves most. It also points to volatility screening as a sensible deployment criterion. Checking the pattern on a wider universe is our immediate next step.

The Ridge-plus-classifier ensemble reaches 53.5\% pooled, below Ridge alone. This one surprised us: the ensemble had beaten Ridge on an earlier one-year window. The reversal is what pushed us to the five-year protocol used throughout this section, and it is why Ridge stays as the production model. A model trained jointly across all five tickers managed 51.0\% and was also set aside, which suggests per-ticker specialisation carries real information at this scale.

%% [MEANING] paraphrase ended mid-sentence ("and not yet a present.")
We also tried gradient boosting on a chronological tune-test split, choosing hyperparameters on the first 70\% of each ticker's history and scoring once on the untouched 30\%. It reached 59.2\% against Ridge's 55.3\%, ahead on four of five tickers. Since we compared four model variants, the appropriate correction leaves that margin short of conventional significance ($p=0.22$), so we treat gradient boosting as the most promising avenue for the next iteration and not as a present claim. Settling it needs roughly 1{,}290 walk-forward steps against the 600 we have, which the expanded universe in Section~\ref{sec:power} would supply.

\subsection{Risk Manager Integration}

%% [MEANING] the paraphrased execution-lag sentence was circular and unreadable
Table~\ref{tab:risk} isolates the Risk Manager's position-sizing rule. To separate that contribution from forecast quality, the driving signal is deliberately a 5/20-day moving-average crossover rather than the Analyst Agent, so the table measures the sizing policy under a signal of known and unremarkable quality. Positions are held at the profile's fixed capital fraction, long on a bullish crossover and short on a bearish one, and the signal observed at one session's close is applied to the next session's return.

\begin{table}[!ht]
\centering
\setlength{\belowcaptionskip}{3pt}
\caption{Risk Manager position sizing vs.\ buy-and-hold under a fixed 5/20 moving-average signal (medium profile, 5-year, 5 tickers). Figures describe this window; see Section~\ref{sec:power}.}
\vspace{-2.5mm}
\begin{tabular}{
  >{\arraybackslash}m{25pt}
  >{\centering\arraybackslash}m{55pt}
  >{\centering\arraybackslash}m{25pt}
  >{\centering\arraybackslash}m{55pt}
  >{\centering\arraybackslash}m{25pt}
  }
  \toprule
  \textbf{Ticker} & \textbf{Model Ret. (\%)} & \textbf{Sharpe} & \textbf{B\&H Ret. (\%)} & \textbf{Sharpe}\\
  \midrule
    AAPL & -1.00 & -0.44 & 17.56 & 0.72 \\
    MSFT & -1.11 & -0.55 & 7.14  & 0.39 \\
    AMZN & -0.79 & -0.26 & 7.41  & 0.38 \\
    TSLA & 3.30  & 0.64  & 13.00 & 0.50 \\
    NVDA & 2.31  & 0.52  & 63.93 & 1.21 \\
  \bottomrule
\end{tabular}
\label{tab:risk}
\end{table}

%% [MEANING] "If there were exposure on the return side..." reversed the sense
The effect on tail risk is large and uniform. Maximum drawdown falls on every ticker we tested: from $-33.4\%$ to $-7.6\%$ on AAPL, and from $-66.4\%$ to $-8.9\%$ on NVDA. That is a four- to sevenfold reduction, and the sign holds across the whole universe. It is the strongest and most consistent result we obtained. The trade-off is exposure. Capital goes to work only on an active signal and is capped by profile, so the policy behaves as a capital-preservation overlay and not as a return maximiser. Lifting exposure through sustained uptrends with a regime-aware sizing rule would address the return side directly.

Because the driving signal is a controlled moving-average crossover, the drawdown reduction belongs to the sizing policy itself and should transfer to any signal of similar quality. That is a stronger and more portable claim than one tied to a particular forecaster. Measuring the Analyst Agent end to end needs a window with reproducible contemporaneous news, which the archival work in Section~\ref{sec:power} will supply.

Figure~\ref{fig:tsla} shows what this looks like for TSLA, the most volatile name we tested. The lower panel carries the result: the drawdown envelope stays inside roughly seven percent for the whole period, including the 2022 decline that the unmanaged position takes in full. The upper panel shows what that costs in participation. Read together, the two panels describe the policy exactly: capital preservation, bought at a price you can measure and control.

\begin{figure}[!ht]
\centering
\includegraphics[width=\linewidth]{tsla_table3_consistent.png}
\caption{TSLA under the medium risk profile: growth of one unit of capital (top) and drawdown envelope (bottom) against buy-and-hold, generated by the same backtest as Table~\ref{tab:risk}. The signal observed at one session's close is applied to the next session's return. Drawdown stays within about seven percent where buy-and-hold exceeds seventy, and most of the upside is forgone in exchange.}
\label{fig:tsla}
\end{figure}

\subsection{System Latency}

Earlier measurement placed the complete Version~4 pipeline under 1.3 seconds end to end on a four-vCPU, 8-GB reference server. The Analyst Agent now fits an additional logistic classifier at every walk-forward step, adding a small and as yet unquantified increment; re-measuring under the updated model remains pending.

\subsection{Data Requirements for Claims of This Type}
\label{sec:power}

Alongside the results we report how much data each kind of claim actually needs. This literature rarely states that calculation, and setting it out lets a reader see which of our conclusions are settled and tells us how to design the next evaluation.

For directional accuracy the arithmetic is quick. At $n=120$ steps per ticker the standard error of an accuracy estimate is $\sqrt{0.25/120}=4.56$ percentage points, so single-ticker comparisons resolve gaps of about nine points or more. Pooling to $n=600$ brings the standard error down to 2.04 points. That is what puts the aggregate Ridge result comfortably inside the detectable range, and why we report it as our headline accuracy figure.

Risk-adjusted return asks for considerably more. Taking the variance of a difference between two Sharpe ratios on correlated return series, with the strategy-to-benchmark correlation measured from the backtest itself at about 0.55, separating a Sharpe improvement of $+0.30$ needs on the order of $10^{4}$ trading days per ticker. Five years gives us roughly 1{,}250. We therefore report the Sharpe column of Table~\ref{tab:risk} as a description of this window, and let the drawdown reduction carry the risk-management conclusion, since it is large, consistent in sign on every ticker, and does not hinge on resolving a small difference in means.

This gives us a specification for the next release instead of an open-ended caveat. Reaching the detection threshold takes roughly eight times the evaluated days, which we can get by widening the universe from five names to a sector-diversified fifty and holding the five-year window fixed. The harness already runs off a ticker list, so no code changes are needed. Extending to non-US venues, which the deployed system supports today, follows the same route.

%% ============================================================
\section{Discussion}
\label{sec:discussion}
%% ============================================================

Keeping language-model reasoning out of the analytical path buys auditability, which matters under financial regulation and which a fully generative pipeline gives up. The obvious extension puts a model higher in the stack, as a retrieval-augmented coordinator sitting over deterministic subagents, so the system can reason more widely without losing that guarantee. Three directions come directly out of the evaluation. Widening the universe from five names to a sector-diversified fifty, as Section~\ref{sec:power} sets out, brings risk-adjusted comparisons into detection range and tests whether the volatility-linked accuracy pattern holds up. Swapping a transformer encoder such as FinBERT in for the lexicon scorer is the most direct accuracy gain open to the Researcher Agent. A regime-aware sizing rule that scales exposure up through sustained uptrends would address what the current capital-preservation policy costs in participation. FINDEC is a decision-support tool and not an investment advisor; every recommendation says so. Because the framework is open source, anyone building on it can audit the scoring for themselves.

%% ============================================================
\section{Conclusion}
\label{sec:conclusion}
%% ============================================================

FINDEC pairs lexicon sentiment, a Ridge ensemble, and risk-profile position sizing in a CPU-only microservice. Directional accuracy clears the random-walk baseline at 55.7\% ($p=0.006$), and position sizing reduces maximum drawdown on every ticker tested, by as much as sevenfold. Harness, data, and code are released for reuse.

\bibliographystyle{IEEEtran}
\bibliography{references}

\end{document}
